\section{データを要約する}

  \subsection{問い}

    食事に行く店を決める, Amazon で電化製品を選ぶ, そういう場面を思い出してみてほしい.
    こういうとき, とりあえずレビュー評価を確認したりしないだろうか.
    星の平均が4.2の製品と3.1の製品があれば, まず4.2のほうを見てみるといった具合である.

    ここで私たちが無意識に行っていることは「データの要約」と「比較」だ.
    星4.2というのは, 何十人もの評価を1つの数字にまとめた結果である.
    何十人分ものバラバラな評価をそのまま頭で処理するのは無理だから, 代わりに「4.2」という1つの数字を指標にしている.
    そして, 同様にまとめられた「3.1」という数をみて, これは4.2のほうが良さそうだぞ, と判断しているのだ.
    日常的に行っている「たくさんのデータを1つにまとめる」「まとめた数どうしを比べる」という行動は, もっとも身近なデータ分析の一つだ.

    ただ, このやり方だけでいつもよい判断ができるとは限らない.
    同じ星4.0の製品が二つ並んでいたら, 次は何を見て選ぶことになるだろう.
    一つが星5だったとしても, たった一人からの評価しかなかったら, その星5をそのまま信じるのはためらわれる.
    データをまとめてしまうと, その数字以外の情報はどこかへ行ってしまうのだ.

    データの要約とは一体どんなもので, どんな情報がなくなってしまうのだろうか.


  \subsection{平均だけでは見えないもの}

    データを1つの数にまとめる方法でもっともポピュラーなのは, 平均だろう.
    データの値を全部足し合わせて個数で割るだけで求まり, 計算しやすいのが取り柄である.

    しかし, 平均だけを見ていてもうまく拾えないことがある.
    たとえば, 通勤に使う電車の平均遅延時間が「1日あたり2分」だったとする.
    別の路線も平均2分だったとしたら, どちらに乗っても同じような通勤になるだろうか.
    実は片方は毎朝きっちり2分遅れるだけで, もう片方は99\%の日は定刻通りだが残り1\%で200分止まる, ということがありうる.
    同じ平均2分でも, 中身はまるで違うのだ.
    平均という1つの数からは, データが平均のまわりに集まっているのか, それとも遠くまで散らばっているのかが見えない.
    平均だけを頼りに「どちらの路線も同じようなものだろう」と判断してしまうと, 年に数回, 始業に3時間遅れる路線を引き当てることになる.
    平均には現れない\textbf{散らばり}という情報が, ここでは判断を左右している.

    あるいは平均年収のような例もある.
    日本の平均年収は約460万円といわれている\footnote{%
    国税庁「令和5年分民間給与実態統計調査」による. 1年を通じて勤務した給与所得者の平均給与が460万円. 中央値は同調査では公表されていないが, 給与階級別の分布から400万円前後と読み取れる.}が, 実は「平均以下」の年収の人のほうが多数派だ.
    一部の高収入者が全体の平均を引き上げているからである.
    こういうときに使われるのが, データを順番に並べて真ん中に来る値をとる\textbf{中央値}という別のまとめ方だ.
    中央値はだいたい400万円前後で, 平均よりも低い.
    真ん中の順位の人は, 高収入者がどれだけ稼ごうがそう簡単には動かないからである.
    平均460万円という数だけを見て「自分の年収はそれより下だから真ん中以下なんだな」などと捉えると, 実態とはずれた印象を持つことになる.

    平均だけでは, 真ん中の人がどのあたりにいるのかが見えてこない.

    ということは, データのまとめ方は平均だけではないし, それぞれの方法が別の情報を拾っているらしい.
    そもそも「データを1つの数にまとめる」とはどういう操作なのか, 一度整理してみよう.


  \subsection{理論の最小核}

    星4.2の例でやっていたのは, 何十人もの評価を「4.2」という1つの数に置き換えることだった.
    多数の値を1つの数で代表させる, それが要約という操作の中身である.

    なぜわざわざ要約するのかといえば, 大きく2つの利点があるからだ.
    ひとつは, 別の集まりと比べやすくなること.
    冒頭で平均4.2と平均3.1を比べたが, あれもまさに同じことをやっていた.
    別の集まりにそれぞれ同じやり方で要約をかけて, 出てきた数同士を見比べていた.
    バラバラの評価をそのまま眺めて比較するのは難しいが, 1つの数に揃えれば見比べるのはたやすい.

    もうひとつは, その数を「対象そのものの値」として扱えるようになることだ.
    平均4.2の製品があったとき, 私たちはそれを「4.2の価値がある製品」のように受け取る.
    何十人もの評価をまとめた結果にすぎない数字が, いつのまにかその製品自体の値として扱われている.
    こうした置き換えは, 後の章で「推定」という言葉でもう少し丁寧に扱う.
    このように, 要約という操作のおかげで比較や推定が成り立っている.
    一方で前の節で見たように, 要約には情報の欠落がついて回る.
    平均は散らばりを拾わないし, 中央値は全体の合計を語らない.
    どんな要約のしかたも, 必ず何かの情報を犠牲にしている.

    データの要約のしかたはいろいろあるが, 「観測値からひとつの数を計算する」という形そのものは共通している.
    ここで, その数に名前を付けておこう.

    \begin{defi}{統計量}{stat}
      $n$個の観測値 $x_1, x_2, \ldots, x_n$ から計算される数を\textbf{統計量}と呼ぶ.
    \end{defi}

    「観測値から計算される数」というのは, 観測値が決まればおのずと値が定まる, ということだ.
    平均にせよ中央値にせよ, 同じ観測値を与えれば計算結果は同じ数になるので, いずれも統計量だ.
    だから統計量というのは特定の「数」というより, 「観測値を入れると数が出てくるなんらかの手続き」というイメージに近い.

    「平均」「中央値」というのも, 個別の値そのものではなく, それぞれの手続きの名前なのだ.
    どんな統計量を使うかは, 何を読み取って何を捨てるかの選び方でもある.

    年収の例では, 高収入者に引きずられる平均よりも, 真ん中の人の位置を拾う中央値のほうが知りたいことに合っていた.
    知りたいことが変われば, 選ぶべき統計量も変わる.

    統計量は平均や中央値だけではなく, 他にもいろいろある.
    過去の統計学者たちが目的に合わせて考え出してきたもので, ふつうはその中から選んで使う.
    自分で一から考えてもかまわないが, だいたいは先人の知恵で間に合う.
    この章ではまず, 代表的なものから順に見ていく.


  \subsection{手法}

    \subsubsection*{平均}

    最初は平均だ.
    データの値を全部足して個数で割る, という計算はここまで何度も登場した.
    式の形にしておこう.

    \begin{defi}{平均}{mean}
      $n$個の観測値 $x_1, x_2, \ldots, x_n$ に対し,
      \[
        \bar{x} = \frac{1}{n} \sum_{i=1}^{n} x_i
      \]
      を\textbf{平均}（算術平均）と呼ぶ.
    \end{defi}

    式の右辺を読み解いておく.
    $\sum_{i=1}^{n} x_i$ ですべての観測値を足し上げ, それを観測値の個数 $n$ で割っている.
    全員に同じ重みを与えて合計を人数でならした量だ.

    平均はすべての観測値を使うので, データの情報をまんべんなく拾う.
    ただしすべてを等しく扱うぶん, 極端に大きい値や小さい値が1つあるだけで結果が大きく引きずられる.
    年収の例で見たとおりである.

    ところで, この平均という統計量はなぜこれほどポピュラーなのだろうか.
    ただの「全部足して個数で割った数」というだけなら, 他の統計量と並んで一つの選択肢にすぎないはずだ.
    それでも平均が筆頭に挙げられるのは, 確率の言葉で見たときに\textbf{期待値}と一致するからだ.
    期待値とは, データの集まりからランダムに1つ選んだとき, だいたいこの値になる, という意味の量である.
    データを見て「だいたいどれくらいか」と知りたくなるのは自然な問いで, 平均はこの問いに正面から答えている.
    平均と期待値が一致するという事実は, この節の最後に命題として確認する.

    \subsubsection*{中央値と四分位数}

    次は, 年収の例に出てきた中央値だ.
    平均がすべての値を足し合わせるのに対し, 中央値は値を並べて真ん中を取る.

    \begin{defi}{中央値}{median}
      観測値を小さい順に並べたとき, ちょうど真ん中に来る値を\textbf{中央値}と呼ぶ.
      $n$ が奇数なら $\frac{n+1}{2}$ 番目の値,
      $n$ が偶数なら $\frac{n}{2}$ 番目と $\frac{n}{2}+1$ 番目の値の平均とする.
    \end{defi}

    中央値はデータを半分に分ける位置にある値だ.
    極端な値がいくつ混じっても真ん中の順位は動きにくいので, 平均に比べて極端な値に引きずられにくい.
    年収のように一部の値が大きく偏っているデータでは, 中央値のほうが「多くの人の実感」に近い数を返すことが多い.

    同じ発想で, 真ん中以外の位置も取ってみる.
    よく使われるのは, 下から4分の1の位置と4分の3の位置だ.

    \begin{defi}{四分位数と四分位範囲}{iqr}
      観測値を小さい順に並べたとき, 下から $25\%$ の位置にある値を\textbf{第1四分位数} $Q_1$,
      下から $50\%$ の位置にある値を\textbf{第2四分位数} $Q_2$（中央値と一致する）,
      下から $75\%$ の位置にある値を\textbf{第3四分位数} $Q_3$ と呼ぶ.
      $Q_1, Q_2, Q_3$ をまとめて\textbf{四分位数}と呼ぶ.
      $Q_1$ と $Q_3$ の差
      \[
        \mathrm{IQR} = Q_3 - Q_1
      \]
      を\textbf{四分位範囲}（interquartile range）と呼ぶ.
    \end{defi}

    $Q_1$ と $Q_3$ で挟まれた範囲には, ちょうど真ん中の $50\%$ の観測値が入っている.
    四分位範囲 $\mathrm{IQR}$ は, その範囲の幅を1つの数にしたものだ.
    真ん中の半分が狭いところに固まっているのか, 広く散らばっているのかを表している.

    \subsubsection*{最頻値}

    位置の取り方には, 真ん中を狙う以外の考え方もある.
    データの中で一番多く現れている値に注目するやり方だ.

    \begin{defi}{最頻値}{mode}
      観測値の中で最も頻繁に出現する値を\textbf{最頻値}と呼ぶ.
      最頻値は複数存在することもある.
    \end{defi}

    平均や中央値が「全体の真ん中をどう取るか」を考えるのに対し, 最頻値は「どの値が一番多く出ているか」という別の問いに答えている.
    レビュー評価でいえば, 一番票を集めた星の数にあたる.
    データの形と直結した統計量で, 第3章でデータの形を図にするときにまた出てくる.

    \subsubsection*{分散と標準偏差}

    平均も中央値も最頻値も, データのおおよその位置を1つの数にしたものだった.
    しかし電車の遅延の例で判断を左右していたのは, 位置ではなく散らばりのほうだ.
    データが平均のまわりに固まっているのか, 遠くまでばらついているのか.
    今度はこの散らばりを1つの数にしてみたい.

    素朴に考えれば, 各観測値と平均との差（偏差）を合計すればよさそうに思える.
    しかし偏差の合計は必ずゼロになる.
    平均より大きい値と小さい値が打ち消し合うからだ.
    そこで偏差を二乗してから合計する.

    \begin{defi}{分散}{variance}
      $n$個の観測値 $x_1, x_2, \ldots, x_n$ とその平均 $\bar{x}$ に対し,
      \[
        s^2 = \frac{1}{n} \sum_{i=1}^{n} (x_i - \bar{x})^2
      \]
      を\textbf{分散}と呼ぶ\footnote{%
        分散の定義には $n$ で割るものと $n-1$ で割るものがある.
        後者は\textbf{不偏分散}と呼ばれ, 母集団の分散を推定する際に用いられる.
        本章では手元のデータを記述するために $n$ で割る版を $s^2$ と書く.
        推論で用いる $n-1$ で割る版は $u^2$ と書き, 同じ記号を使い回さない.
        $n-1$ で割る理由については後の章で改めて扱う.}.
    \end{defi}

    \begin{defi}{標準偏差}{stddev}
      分散の正の平方根
      \[
        s = \sqrt{s^2}
      \]
      を\textbf{標準偏差}と呼ぶ.
    \end{defi}

    定義式を読み解いておく.
    $(x_i - \bar{x})$ は各観測値の平均からのずれ（偏差）だ.
    偏差を二乗して足し上げ, $n$ で割っている.
    二乗するのは, 正のずれと負のずれが打ち消し合わないようにするためだ.
    $n$ で割るのは, 偏差の二乗の平均を取るためである.
    合計のままだと観測値が多いほど値が膨らんでしまい, 個数の違う集まりどうしで散らばりを比べられない.
    ところで, 打ち消しを防ぐだけなら絶対値を取る手もある.
    それでも二乗の方が標準になっているのは, 数学的に扱いやすい性質（たとえば微分ができる）を多く持つからだ.
    どちらでなければならないと自然に決まるものではなく, 人間の都合で定着した約束である.

    もう一つ, 偏差をまとめて眺める見方がある.
    $n$個の偏差 $(x_1 - \bar{x}), \ldots, (x_n - \bar{x})$ をひとまとめにして, 一本の矢印と考える. このように数の組を一本の矢印として扱ったものをベクトルと呼ぶ.
    偏差が2つなら, 2つの偏差を直角三角形の縦と横に取ったとき, この矢印が斜辺になる.
    三平方の定理から斜辺の長さの二乗は2つの偏差の二乗和であり, 偏差が$n$個でも同じ形が続くため, この矢印の長さの二乗は $\sum_{i=1}^{n}(x_i - \bar{x})^2$ となる.
    したがって標準偏差は, データが平均からどれだけ離れているかを, 矢印の長さを $\sqrt{n}$ でならした距離として読める.
    これは二乗でなければならないと示すものではなく, 二乗を使う定義に別の読み方を与える見方である.

    分散は偏差の二乗の平均なので, 単位が元のデータの二乗になる.
    たとえばデータの単位が「円」なら分散の単位は「円${}^2$」だ.
    標準偏差はその平方根を取ることで, 元のデータと同じ単位に戻す.
    散らばりの大きさを実感するには, 標準偏差のほうが扱いやすい.

    分散と標準偏差は平均と同じく, すべての観測値を使って計算する.
    そのため極端な値の影響を受けやすい.
    偏差を二乗するぶん, その影響は平均よりむしろ増幅される.

    \subsubsection*{他にもいろいろある}

    統計量には, 中心や散らばりを測るもののほかに, 比較や関係を測るものもたくさんある.
    たとえば後の章では, 観測されたデータと期待されるデータのずれを測るカイ二乗統計量や, 2つの量がどれくらい一緒に動くかを表す相関係数が出てくる.
    どれも, それぞれの目的のために考え出されたものだ.
    具体的な中身は, それぞれが必要になる章で紹介する.

    \subsubsection*{平均と期待値の関係}

    平均の項で予告した, 期待値との関係をここで確かめておく.

    \begin{prop}{平均と期待値}{mean-expectation}
      重複を許して並べた $n$ 個の値 $x_1, x_2, \ldots, x_n$ からなる母集団から1つをランダムに取り出すとき, その値を確率変数 $X$ とする.
      このとき $X$ の期待値は
      \[
        E[X] = \frac{1}{n} \sum_{i=1}^{n} x_i = \bar{x}
      \]
      に等しい.
    \end{prop}

    証明はこの節の最後のコラムに置く.

    \subsubsection*{統計量を組み合わせて見る}

    ここまで導入した統計量を, ひとつのデータに対して並べて使ってみよう.

    あるECサイトに2つの商品 A, B があり, どちらも20件のレビューがついている.
    平均評価はどちらもちょうど 4.0 だ.
    平均だけ見れば, この2つは同じくらいの「良い商品」に映る.

    ところが評価の中身を見ると, 商品 A は 4 点が10件, その両側の 3 点と 5 点が 5 件ずつで, 4 点を中心におおむね固まっている. 一方の商品 B は 5 点が15件と 1 点が 5 件で, 中間がまったくない.
    熱烈な支持と強い不満が同居している商品だ.

    平均以外の統計量を並べてみる.

    \begin{center}
      \begin{tabular}{l|cc}
        & 商品 A & 商品 B \\ \hline
        平均 & 4.0 & 4.0 \\
        中央値 & 4 & 5 \\
        標準偏差 & 0.71 & 1.73 \\
      \end{tabular}
    \end{center}

    平均は同じだが, 中央値と標準偏差が大きく違う.
    中央値が 5 ということは, 並べた20件の真ん中の人は満点をつけている.
    それでも平均が 4.0 まで下がっているのは, 1 点の評価が全体を引きずり下ろしているからだ.
    標準偏差は商品 A の倍以上で, 評価のばらつきの大きさが数字にもはっきり現れている.

    平均1つでは見えなかったことが, 統計量を組み合わせて並べると見えてくる.
    どの統計量にも単独では死角があるから, 角度の違う物差しを何本か当てて, 互いの死角を補う.
    平均が同じだから同じ商品, とはならないのだ.

    \begin{mybox}{平均と期待値の一致の証明}
      母集団の中で値 $a$ が $k_a$ 回現れているとする.
      1つをランダムに取り出すとき, 値 $a$ が出る確率は $k_a / n$ である.
      期待値の定義から,
      \[
        E[X] = \sum_{a} a \cdot \frac{k_a}{n} = \frac{1}{n} \sum_{a} a \cdot k_a.
      \]
      ここで $\sum_{a} a \cdot k_a$ は, 各値 $a$ をその出現回数 $k_a$ ぶん足し合わせたものなので, 元の $x_1 + x_2 + \cdots + x_n$ に等しい.
      したがって
      \[
        E[X] = \frac{1}{n} \sum_{i=1}^{n} x_i = \bar{x}.
      \]
    \end{mybox}


  \subsection{実践}

    この章で見てきた統計量を, 本書を通して扱うランニング例のデータに当てはめてみる.

    扱うのは, あるサブスク型動画配信サービスから抽出された邦画(実写)作品80本の標本データだ.
    各作品について, 公開後30日間の視聴回数, 製作費, 平均評価, 評価件数, ジャンル, 原作の有無などが記録されている.
    まずは分析の中心となる\textbf{公開後30日間の視聴回数}に注目し, この80本の視聴回数の集まりがどんな姿をしているかを統計量で要約してみよう.

    \subsubsection*{平均と中央値を計算する}

    まず平均を計算してみると, 80本の視聴回数の平均はおよそ $209.4$ 万回であった.
    一方, 中央値は $107.1$ 万回である.

    平均と中央値が大きく食い違っている.
    平均は中央値の約2.0倍だ.

    視聴回数の集まりには, 大ヒット作品が少数混じっている.
    その視聴回数は桁違いに大きく, 平均を引きずり上げている.
    年収の例で見たのと同じ構造である.
    視聴回数のように一部の値が突出するデータでは, 平均は「ほとんどの作品の視聴回数」よりかなり大きくなる.
    「典型的な作品はどのくらい見られているか」を知りたいのなら, 平均より中央値のほうが実感に近い数を返す.

    \subsubsection*{散らばりを測る}

    次に散らばりを見る.
    標準偏差を計算するとおよそ $239.3$ 万回となる.
    平均とほぼ同じくらいの値だ.
    散らばりの幅が平均そのものと同じくらいあるのだから, 作品ごとの視聴回数は相当大きくばらついている.

    四分位範囲も計算しておこう.
    第1四分位数 $Q_1$ は $57.2$ 万回, 第3四分位数 $Q_3$ は $263.4$ 万回で, 差を取ると $\mathrm{IQR}$ は $206.2$ 万回になる.
    つまり, 視聴回数を小さい順に並べたときに真ん中50\%にあたる作品たちは, この幅に収まっている.

    標準偏差は全作品を使った散らばりの指標で, 突出した作品の影響を強く受ける.
    四分位範囲は両端の25\%ずつを切り捨てたうえでの散らばりなので, 突出した作品の影響を受けにくい.
    どちらも散らばりを測る統計量だが, 答えている問いが少しずつ違う.

    \subsubsection*{並べてみる}

    ここまでに計算した統計量を一つの表にまとめてみる.

    \begin{center}
      \begin{tabular}{l|c}
        統計量 & 視聴回数 (万回) \\ \hline
        平均 & 209.4 \\
        中央値 & 107.1 \\
        標準偏差 & 239.3 \\
        $Q_1$ & 57.2 \\
        $Q_3$ & 263.4 \\
        $\mathrm{IQR}$ & 206.2 \\
      \end{tabular}
    \end{center}

    平均は中央値から大きく離れ, 標準偏差は平均と同じくらい大きく, 四分位範囲はその標準偏差より狭い.
    この3つを並べると, 80本の視聴回数の集まりは, 多くの作品が中央値あたりに固まっている一方で, 一部の作品が大きく上に飛び出している, という姿が浮かんでくる.

    平均だけ見ていては気づけなかった集まりの形が, 統計量を組み合わせることで見えてきた.


  \subsection{結果と次の一手}

    この章では, データの集まりを統計量で捉えるという発想を導入し, 平均, 中央値, 四分位数, 最頻値, 分散, 標準偏差といった基本的な物差しを見てきた.
    どの統計量も観測値から計算される数だが, 答えている問いと, 落としている情報が違う.
    複数の統計量を組み合わせて見ることで, 単独では見えなかった集まりの性格が浮かび上がってくる.

    ただし, どの統計量も, データの一面を1つの数に切り取ったものであることに変わりはない.
    切り取るかぎり, データ全体がどんな姿をしているかという情報はどうしても落ちる.
    その姿を見るには, 数を図に置き換えるという別の方法が要る.
    それが次の章で扱う\textbf{可視化}だ.

  \begin{mybox}{ケトレーと「平均人」}
    19世紀のベルギーの統計学者アドルフ・ケトレーは, 人体の寸法や犯罪率を平均で要約することに情熱を注ぎ, 「平均人（l'homme moyen）」という概念を打ち出した.\footnotemark
    社会の典型を平均で語ろうとする発想だ.

    しかしケトレーは平均にこだわるあまり, 平均から外れた個人を「ずれ」「異常」と見なす方向に進んでしまった.
    平均1.6人の子供を持つ世帯が存在しないのと同じで, 「平均人」もまた, 現実には誰も該当しない値だ.
    平均を集まりの代表として使うことと, 平均を「あるべき姿」として個人に押しつけることの間には, 大きな隔たりがある.

    統計量は集まりを測る物差しであって, 個人に当てる物差しではない.
    平均を扱うときに, 一度は思い出しておきたい区別だ.
  \end{mybox}
  \footnotetext{%
  Quetelet, A. (1835). \textit{Sur l'homme et le d\'eveloppement de ses facult\'es, ou Essai de physique sociale}. Paris: Bachelier.}
